%% F12 --- Reguły różniczkowania i reguła łańcuchowa
%%
%% Zalążek. Poniższy opis jest kontraktem tego programu: co ma uzasadnić i co
%% ma dać czytelnikowi. Napisanie programu oznacza usunięcie bloku
%% \programstub{} i zastąpienie go programem. Jeśli po przerobieniu materiału
%% opis okaże się błędny, popraw go i odnotuj sprzeczność w CLAUDE.md w sekcji
%% *Resolved questions*.

\program{Reguły różniczkowania i reguła łańcuchowa}
\label{prog:F12}

\programstub{%
\textit{[Opis do przetłumaczenia --- patrz wydanie angielskie.]}

Argument: four rules and one composition rule cover everything the book will
differentiate. Payoff: backpropagation is the chain rule applied to a
composition of layers, and that sentence is the hinge of the entire book.
The longest Foundation program for that reason.

\medskip
\textbf{Planowana długość:} 55 ramek.
}
